\def\SigleNumero{STT-1000}
\def\NomCours{Probabilités et statistique}
\def\DateEvaluation{27 février 2025}
\def\HeureEvaluation{18h30 à 21h20}
\def\Session{}
\def\NomEvaluation{Examen 1}
\def\Ponderation{40}
% Ne rien mettre entre les accolades si l'enseignant (ou l'enseignante) est aussi le coordonnateur (ou la coordonnatrice)
\def\Coordonnatrice{}
\def\Coordonnateur{}

% Ne rien mettre entre les accolades s'il y a plus d'une section
\def\Enseignant{}
\def\Enseignante{}

% Ne rien mettre entre les accolades s'il s'agit d'un cours à section unique
% Mettre le nom de chacun des enseignants et enseignantes entre les accolades
% Une case à cocher sera générée pour chacune des sections
\def\SectionA{Rachid Kandri-Rody}
\def\SectionB{Julien Miron}
\def\SectionC{}
\def\SectionD{}


\def\Directives{
\begin{enumerate}
\item Matériel autorisé:
\begin{itemize}
\item Calculatrice scientifique {\bf autorisée par la faculté des sciences et de génie};
\item Le feuillet de formules fourni avec l'examen. Celui-ci inclut la table de la loi normale. Veuillez ne rien écrire sur le feuillet de formules.
\end{itemize}
\item Assurez-vous que cette évaluation comporte \numquestions ~exercices répartis sur \numpages{}~pages, pour un total de \numpoints{}~points.\
\item \textbf{Veuillez ne pas dégrafer votre examen.}
\item Ajoutez vos initiales dans le bas de chaque page.
\item Vous devez\textbf{ justifier toutes vos réponses} par une démarche claire, sauf pour les questions à choix de réponse.\
\item Pour les vrais ou faux, vous pouvez mettre un X, un crochet ($\checkmark$) ou remplir la case.\
\item Les vrais ou faux demande une justification. Aucun point ne sera donné sans justification, même si vous cochez la bonne réponse.
\item Le verso des feuilles peut être utilisé comme brouillon, \textbf{mais ne sera pas corrigé}. Attention à ne rien dégrafer.
\end{enumerate}
}



\documentclass{Evaluation}
\usepackage{multicol,graphicx}
\begin{document}

\newpage
%%%%%%%%%%%%%%%% QUESTION %%%%%%%%%%%%%%%%
\begin{questions}

%11111 Rachid: Proba cdn et Bayes
\question
La tuberculose bovine est une maladie très contagieuse. Cette maladie est apparue dans  le cheptel bovin d'un pays et touche $10 \%$ de ce cheptel. Des études statistiques ont montré que la probabilité pour un animal d'avoir un test positif à cette maladie sachant qu'il est malade est de $95 \%$  et que celle d'avoir un test négatif sachant qu'il n'est pas atteint par la maladie est de $90 \%$ . Définissons les événements suivants :\\
\begin{tabular}{ l}
$T$ : avoir un test positif à cette maladie pour un bovin choisi au hasard parmi le cheptel.\\
$M$ : être malade pour un bovin choisi au hasard parmi le cheptel.
\end{tabular}

\medskip
\begin{parts}
\part [4] Déterminez la probabilité d'avoir un test positif à cette maladie pour un bovin choisi au hasard parmi le cheptel sachant que celui-ci n'est pas malade.

\vspace{7cm}
\part [4] Calculez la probabilité pour un bovin choisi au hasard parmi le cheptel d'avoir un test positif à cette maladie.

\vspace{7cm}
\part [4] Calculez la probabilité qu'un bovin choisi au hasard parmi le cheptel soit malade sachant que le test est positif.
\end{parts}


\newpage
%22222 Rachid:  V.a discrete
\question 
Soit $X$ une variable aléatoire ayant une espérance mathématique \textbf{nulle} et dont la fonction de masse est donnée par
$$
\begin{array}{c|ccc}
x&-1&0&1\\\hline
&&&\\[-0.2cm]
p(x) &\;\;\;\;a\;\;\;\;&\;\;\;\;b\;\;\;\;&\;\;\;\;2\,b\;\;\;\;
\end{array}
$$
\begin{parts}
\part [4] Montrez que $\displaystyle a=\frac{2}{5}$ et $\displaystyle b=\frac{1}{5}$.

\vspace{6cm}
\part [4] Calculez la variance de $X$.

\vspace{6cm}

\part [4] Calculez $E(-3X+1)$ et $V(-3X+1)$.
\end{parts}


\newpage
%33333 Rachid: Lois binomiale, géométrique et Pascal
\question
Un vendeur d’ordinateurs a constaté que 30\% des clients qui entrent dans son magasin, faisaient l’achat d’un ordinateur et que les clients achètent les ordinateurs indépendamment les uns des autres.
Définissez clairement les variables aléatoires qui interviennent dans vos calculs pour répondre aux questions suivantes.
\begin{parts}
\part [4] Quatre clients sont dans le magasin. Quelle est  la probabilité qu’au moins un d’entre eux achète un ordinateur?  

\vspace{7cm}
\part [4] Calculez la probabilité que 4 clients doivent entrer dans le magasin pour que le magasin vende son premier ordinateur? 

\vspace{7cm}

\part [4] Combien de clients, en moyenne, doivent entrer dans le magasin pour que le magasin vende son troisième ordinateur?  
\end{parts}



\newpage
%44444 Rachid: Loi normale
\question
Le diamètre $X$ (exprimé en cm) des tomates livrées à une usine d'emballage
américaine suit une loi normale $\mathcal{N}(7,\ \sigma^2)$ où l'écart-type $\sigma$ de $X$ est inconnu. Un tri automatique rejette
toutes les tomates dont le diamètre n'est pas compris entre 6 cm et 8 cm.
\begin{parts}
\part [6] On constate que 10\% des tomates livrées sont rejetées par ce procédé de tri.\\ Calculez l'écart type $\sigma$.

\vspace{10cm}
\part [6] Le  directeur  veut  réduire  à  5\%  le  pourcentage  de  tomates  rejetées  lors  du  tri.  Ne
pouvant  agir  sur  les  livraisons,  il  installe  un  système  de  tri  qui  rejette  les  tomates  de
diamètre inférieur à $7-s$ ou supérieur à $7+s$. \\
Calculez s.
\end{parts}


\newpage
%55555 Rachid: Lois de Poisson, ex^ponentielle et Gamma.
\question
Les avions d’une grande compagnie aérienne atterrissent, dans un aéroport très achalandé, selon une loi de Poisson avec une moyenne de six avions à l’heure.

Définir clairement les variables aléatoires qui interviennent dans vos calculs pour répondre aux questions suivantes.
\begin{parts}
\part [4] Calculez la probabilité de voir atterrir 2 avions ou plus en 15 minutes.

\vspace{7cm}
\part [4] Calculez la probabilité de devoir attendre plus de 20 minutes pour voir atterrir un avion. 

\vspace{7cm}
\part [4] Combien de temps, en moyenne, devriez-vous attendre pour voir atterrir, dans cet aéroport,  le cinquième avion de cette compagnie.
\end{parts}



\newpage
%66666 Julien
\question  Soit la variable aléatoire $X$ qui suit une loi uniforme sur l'intervalle $[a,b]$. Nous savons que $E(X)=4$ et $V(X)=3$. 

\begin{parts}
    \part[5] Donnez la fonction de densité de $X$.
\begin{solution}
    Il faut déduire $a$ et $b$ à partir de l'espérance et de la variance.

    Nous avons 
    \begin{align*}
        E(X)=&\frac{a+b}{2}\\
        4=&\frac{a+b}{2}\\
        \text{et} &\\
        V(X)=&\frac{(b-a)^2}{12}\\
        3=&\frac{(b-a)^2}{12}
    \end{align*}

    On déduit que $a=8-b$, donc $36 = (b-(8-b))^2=(2b-8)^2$. Il existe donc deux solutions pour $b$. 
    \begin{align*}
        6=2b-8 & & -6=2b-8\\
        b=7 & & b =1
    \end{align*}

Si $b=7$, nous avons $a=1$ et si $b=1$, nous avons $a=7$, ce qui est contradictoire avec l'intervalle $ [a,b]$, où il est sous-entendu que $b>a$.

La fonction de densité est donc

$$f(x)=\left\{
			\begin{array}{ll}
   \displaystyle \frac{1}{6}  & \mbox{si } x\in [1,7] \\
    \displaystyle 0  & \mbox{ sinon, }
			\end{array}
			\right.$$

\end{solution}
    \vspace{10cm}
    \part[4] Sachant que la fonction de répartition de $X$ est donnée par $F(x)=\displaystyle \frac{x-1}{6}$ pour $x\in [a,b]$ et $F(x)=0$ sinon, calculez $P(X>3|X>2)$.
    \begin{solution}
        \begin{align*}
            P(X>3|X>2)=&\frac{P(X>3 \cap X>2)}{P(X>2)}\\
            =& \frac{P(X>3)}{P(X>2)}\\
            =& \frac{1-F(3)}{1-F(2)}\\
            =& \frac{1-\frac{3-1}{6}}{1-\frac{2-1}{6}}\\
            =& \frac{2/3}{5/6}\\
            =& \frac{4}{5}
        \end{align*}
    \end{solution}
\end{parts}
\newpage
%77777 Julien
\question  On considère deux variables aléatoires discrètes $X_1$ et $X_2$ pour lesquelles la fonction de probabilité de masse conjointe est donnée par 

\begin{center}
\begin{tabular}{|cc||c|c|c|}\hline
$p(x_1,x_2)$ & & \multicolumn{3}{c|}{$x_2$}  \\ \cline{3-5}
 & & 0 & 2 & 4  \\ \hline
 & 0 & 0 & 3/15 & 3/15  \\ \hline
$x_1$ & 1 & 2/15 & 6/15  & 0  \\ \hline
 & 2 & 0 &  1/15  &  0 \\ \hline
\end{tabular}
\end{center}

\begin{parts}
\part[4] Donnez la loi conditionnelle de $X_2$ sachant que $X_1 = 1$. Justifiez vos calculs.

\begin{solution}
 Par définition de la probabilité conditionnelle, on a 
$$P(X_2 = x | X_1 = 1) = \frac{P(X_2=x \cap X_1=1)}{P(X_1=1)}$$

On calcule d'abord $$P(X_1=1) = P(X_1=1 \cap X_2 = 0) + P(X_1=1 \cap X_2 = 2) + P(X_1=1 \cap X_2 = 4)  = 2/16 + 6/15 + 0 = 8/15.$$

On a donc 
$$P(X_2 = 0 | X_1 = 1) = \frac{P(X_2=0 \cap X_1=1)}{P(X_1=1)} = \frac{2/15}{8/15} = 1/4$$

On obtient de la même façon $P(X_2 = 2 | X_1 = 1) = \frac{6/15}{8/15} = 3/4$ et $P(X_2 = 2 | X_1 = 1) = \frac{0}{8/15} = 0.$
   
\end{solution}


\vspace{5cm}

\part[4] Sachant que $E(X_1) = 2/3$ et que $E(X_2) = 2$, calculez la covariance entre $X_1$ et $X_2$. Laissez des traces de votre démarche.

\begin{solution}
Il n'y a que deux combinaisons des variables $X_1$ et $X_2$ qui donnent un produit non-nul. On a donc
$$E(X_1X_2) = 2*6/15 + 4*1/15 = 16/15$$
La covariance est donc $E(X_1X_2)-E(X_1)E(X_2) = 16/15 - 2/3*2 = 16/15-4/3 = -4/15.$
\end{solution}

\vspace{6cm}

\part[4] Montrez, sans utiliser le résultat obtenu en (b), que les variables aléatoires $X_1$ et $X_2$ ne sont pas indépendantes.

\begin{solution}
On a calculé en (a) que $P(X_1=1) = 8/15$. 
On voit aussi que $P(X_2=4) = 3/15. $
Si Les variables aléatoires étaient indépendantes, alors on aurait $P(X_1=1 \cap X_2 = 4) = P(X_1=1)P(X_2=4) = 8/15 * 3/15 = 24/225.$ Or cette probabilité conjointe vaut plutôt zéro, montrant que les variables aléatoires ne sont pas indépendantes. 

(Plusieurs autres réponses possibles.)
\end{solution}

\end{parts}
\newpage
%88888 Julien
\question Soit $A$, $B$ et $C$ trois événements avec $P(A)=0,5$, $P(B)=0,5$, $P(C)=0,2$, $P(A\cap B)=0,25$. Nous avons aussi que $C$ est indépendant de l'événement $A$ et de l'événement $B$. Dites si chacun des énoncés ci-dessous est vrai ou faux. Votre réponse doit être justifiée en mots ou par des calculs. Un vrai ou faux sans justification n'aura aucun point.

\vspace{0.5cm}
\begin{parts}

\part[3] Les événements $A$ et $B$ sont indépendants.

     \begin{oneparcheckboxes}
		\correctchoice VRAI
		\choice FAUX
\end{oneparcheckboxes}\\
{\bf Justification:}
\begin{solution}
    Comme $P(A)P(B)=0,5 \times 0,5 = 0,25 = P(A\cap B)$, les événements sont indépendants.
\end{solution}
\vspace{4cm}

\part[3] $P(A\cup C)= 0,6$.

\begin{oneparcheckboxes}
		\correctchoice VRAI
		\choice FAUX
\end{oneparcheckboxes}\\
{\bf Justification:}
\begin{solution}
    Comme $A$ et $C$ sont indépendants, $P(A\cap C) = 0,5 \times 0,2 =0.1$.

    Donc $P(A \cup C)=P(A)+P(C)-P(A\cap C)=0,5+0,2-0,1=0,6$
\end{solution}
\vspace{4cm}
\part[3] Les événements $A$ et $B$ forment une partition de l'espace échantillonnal.

     \begin{oneparcheckboxes}
		\choice VRAI
		\correctchoice FAUX
\end{oneparcheckboxes}\\
{\bf Justification:}
\begin{solution}
    Pour qu'ils forment une partition de $\mathcal{U}$, ils doivent être mutuellement exclusifs, i.e. $A\cap B= \emptyset$, ce qui devrait donner $P(A\cap B)=0$, ce qui n'est pas le cas.
\end{solution}
\end{parts}

\newpage

\question Après un certain nombre de tours pendant une partie de \textit{Scrabble}, les 15 lettres suivantes n'ont toujours pas été choisies et sont donc encore dans le sac:
$$\{a, e, i, g, h, l, n, o, p, t, t, v, v, w, x\}.$$

\begin{parts}
    \part[4] Si Alex doit piger 4 lettres au hasard dans le sac, quelle est la probabilité que ces lettres soit les lettres de son prénom, dans le bon ordre? Laissez des traces de votre démarche.
\begin{solution}
    Soit l'événement $A:$ Alex pige les lettres de son nom dans l'ordre.
    
    Comme l'ordre est important, la probabilité est

    \begin{align*}
        P(A)&=\frac{1}{15}\times\frac{1}{14}\times\frac{1}{13}\times\frac{1}{12}\\
        P(A)&= \frac{1}{32760}
    \end{align*}
\end{solution}

    \vspace{8cm} 
    \part[6] Supposons qu'Alex a pigé $\{i, g, t, w\}$. Il reste donc 
$\{a, e, h, l, n, o, p, t, v, v, x\}$ dans le sac. C'est maintenant au tour d'Eva de piger 4 lettres. Quelle est la probabilité que parmi les 4 lettres qu'elle doit piger, elle pige les lettres qui forment son prénom (peu importe l'ordre)? Laissez des traces de votre démarche.
\begin{solution}
     Soit l'événement $E:$ Eva pige au moins les 3 lettres de son nom peu importe l'ordre.
    
    Comme l'ordre n'est pas important.

    \begin{align*}
        P(E)&= \frac{\binom{1}{1}\binom{1}{1}\binom{2}{1}\binom{8}{1}}{\binom{11}{4}}\\
        &= \frac{16}{330}\\
        &\approx 0,0485
    \end{align*}
\end{solution}
\end{parts}






\end{questions}
\end{document} 
